\section{Dimension}
\label{section-dimension}

\noindent
The dimension of a scheme is just the dimension of its underlying
topological space.

\begin{definition}
\label{definition-dimension}
Let $X$ be a scheme.
\begin{enumerate}
\item The {\it dimension} of $X$ is just the dimension of $X$
as a topological spaces, see
Topology, Definition \ref{topology-definition-Krull}.
\item For $x \in X$ we denote $\dim_x(X)$ the dimension of the underlying
topological space of $X$ at $x$ as in
Topology, Definition \ref{topology-definition-Krull}.
We say $\dim_x(X)$ is the {\it dimension of $X$ at $x$}.
\end{enumerate}
\end{definition}

\noindent
As a scheme has a sober underlying topological space
(Schemes, Lemma \ref{schemes-lemma-scheme-sober})
we may compute the dimension of $X$ as the supremum of the lengths $n$
of chains
$$
T_0 \subset T_1 \subset \ldots \subset T_n
$$
of irreducible closed subsets of $X$, or as the supremum of the lengths $n$
of chains of specializations
$$
\xi_n \leadsto \xi_{n - 1} \leadsto \ldots \leadsto \xi_0
$$
of points of $X$.

\begin{lemma}
\label{lemma-dimension}
Let $X$ be a scheme. The following are equal
\begin{enumerate}
\item The dimension of $X$.
\item The supremum of the dimensions of the local rings of $X$.
\item The supremum of $\dim_x(X)$ for $x \in X$.
\end{enumerate}
\end{lemma}

\begin{proof}
Note that given a chain of specializations
$$
\xi_n \leadsto \xi_{n - 1} \leadsto \ldots \leadsto \xi_0
$$
of points of $X$ all of the points $\xi_i$ correspond to prime ideals
of the local ring of $X$ at $\xi_0$ by
Schemes, Lemma \ref{schemes-lemma-specialize-points}.
Hence we see that the dimension of $X$ is the supremum of the dimensions
of its local rings. In particular $\dim_x(X) \geq \dim(\mathcal{O}_{X, x})$
as $\dim_x(X)$ is the minimum of the dimensions of open neighbourhoods of
$x$. Thus $\sup_{x \in X} \dim_x(X) \geq \dim(X)$. On the other hand,
it is clear that $\sup_{x \in X} \dim_x(X) \leq \dim(X)$
as $\dim(U) \leq \dim(X)$ for any open subset of $X$.
\end{proof}

\begin{lemma}
\label{lemma-codimension-local-ring}
Let $X$ be a scheme. Let $Y \subset X$ be an irreducible closed
subset. Let $\xi \in Y$ be the generic point. Then
$$
\text{codim}(Y, X) = \dim(\mathcal{O}_{X, \xi})
$$
where the codimension is as defined in
Topology, Definition \ref{topology-definition-codimension}.
\end{lemma}

\begin{proof}
By Topology, Lemma \ref{topology-lemma-codimension-at-generic-point}
we may replace $X$ by an affine open neighbourhood of $\xi$. In this
case the result follows easily from
Algebra, Lemma \ref{algebra-lemma-irreducible-components-containing-x}.
\end{proof}

\begin{lemma}
\label{lemma-generic-point}
Let $X$ be a scheme. Let $x \in X$. Then $x$ is a generic point of
an irreducible component of $X$ if and only if $\dim(\mathcal{O}_{X, x}) = 0$.
\end{lemma}

\begin{proof}
This follows from Lemma \ref{lemma-codimension-local-ring} for example.
\end{proof}

\begin{lemma}
\label{lemma-locally-Noetherian-dimension-0}
A locally Noetherian scheme of dimension $0$ is a disjoint
union of spectra of Artinian local rings.
\end{lemma}

\begin{proof}
A Noetherian ring of dimension $0$ is a finite product of Artinian local
rings, see
Algebra, Proposition \ref{algebra-proposition-dimension-zero-ring}.
Hence an affine open of a locally Noetherian scheme $X$ of dimension
$0$ has discrete underlying topological space. This implies that
the topology on $X$ is discrete. The lemma follows easily from these
remarks.
\end{proof}

\begin{lemma}
\label{lemma-dimension-zero}
\begin{reference}
Email from Ofer Gabber dated June 4, 2016
\end{reference}
Let $X$ be a scheme of dimension zero. The following are equivalent
\begin{enumerate}
\item $X$ is quasi-separated,
\item $X$ is separated,
\item $X$ is Hausdorff,
\item every affine open is closed.
\end{enumerate}
In this case the connected components of $X$ are points and every
quasi-compact open of $X$ is affine. In particular, if $X$
is quasi-compact, then $X$ is affine.
\end{lemma}

\begin{proof}
As the dimension of $X$ is zero, we see that for any affine open
$U \subset X$ the space $U$ is profinite and satisfies
a bunch of other properties which we will use freely below, see
Algebra, Lemma \ref{algebra-lemma-ring-with-only-minimal-primes}.
We choose an affine open covering $X = \bigcup U_i$.

\medskip\noindent
If (4) holds, then $U_i \cap U_j$ is a closed subset of
$U_i$, hence quasi-compact, hence $X$ is quasi-separated,
by Schemes, Lemma \ref{schemes-lemma-characterize-quasi-separated},
hence (1) holds.

\medskip\noindent
If (1) holds, then $U_i \cap U_j$ is a quasi-compact open
of $U_i$ hence closed in $U_i$. Then $U_i \cap U_j \to U_i$
is an open immersion whose image is closed, hence it is a
closed immersion. In particular $U_i \cap U_j$ is affine
and $\mathcal{O}(U_i) \to \mathcal{O}_X(U_i \cap U_j)$ is surjective.
Thus $X$ is separated
by Schemes, Lemma \ref{schemes-lemma-characterize-separated},
hence (2) holds.

\medskip\noindent
Assume (2) and let $x, y \in X$. Say $x \in U_i$. If $y \in U_i$
too, then we can find disjoint open neighbourhoods of $x$ and $y$
because $U_i$ is Hausdorff. Say $y \not \in U_i$ and $y \in U_j$.
Then $y \not \in U_i \cap U_j$ which is an affine open of $U_j$
and hence closed in $U_j$. Thus we can find an open neighbourhood
of $y$ not meeting $U_i$ and we conclude that $X$ is Hausdorff,
hence (3) holds.

\medskip\noindent
Assume (3). Let $U \subset X$ be affine open.
Then $U$ is closed in $X$ by Topology, Lemma
\ref{topology-lemma-quasi-compact-in-Hausdorff}.
This proves (4) holds.

\medskip\noindent
Assume $X$ satisfies the equivalent conditions (1) -- (4).
We prove the final statements of the lemma. Say $x, y \in X$
with $x \not = y$. Since $y$ does not specialize to $x$ we can
choose $U \subset X$ affine open with $x \in U$ and $y \not \in U$.
Then we see that $X = U \amalg (X \setminus U)$ is a decomposistion
into open and closed subsets which shows that $x$ and $y$ do not belong
to the same connected component of $X$. Next, assume $U \subset X$
is a quasi-compact open. Write $U = U_1 \cup \ldots \cup U_n$ as a union
of affine opens. We will prove by induction on $n$ that $U$ is affine.
This immediately reduces us to the case $n = 2$. In this case we have
$U = (U_1 \setminus U_2) \amalg (U_1 \cap U_2) \amalg (U_2 \setminus U_1)$
and the arguments above show that each of the pieces is affine.
\end{proof}

\begin{lemma}
\label{lemma-isolated-dim-0-locally-noetherian}
Let $x$ be a point of a locally Noetherian scheme $X$.
Then $\dim_x(X) = 0$ if and only if $x$ is an isolated point of $X$.
\end{lemma}

\begin{proof}
If $x$ is an isolated point, then $\{x\}$ is an open subset of $X$
(Topology, Definition \ref{topology-definition-isolated-point})
and hence $\dim_x(X) = \dim(\{x\}) = 0$ by definition. Conversely, if
$\dim_x(X) = 0$, then there exists an open neighbourhood $U \subset X$
of $x$ such that $\dim(U) = 0$
(Topology, Definition \ref{topology-definition-Krull}).
By Lemma \ref{lemma-locally-Noetherian-dimension-0}
we see that the topology on $U$ is discrete and hence
$x$ is an isolated point.
\end{proof}





